% Technical report fragment for inclusion in a larger LaTeX document.
% Suggested packages: amsmath, booktabs, graphicx.
% Suggested bibliography keys: 3gpp38331, 3gpp38133, schulman2017ppo.

\section{Technical Description of the Event-Triggered PPO Handover Framework}
\label{sec:project-technical-report}

This section documents the current implementation and experimental status of
the project. Its purpose is to provide a complete technical reference from
which the final paper structure can later be derived. The framework studies
adaptive handover decisions for a cellular-connected unmanned aerial vehicle
(UAV) using pre-generated RSRP and SINR traces. A protocol engine first
reproduces the main timing and state transitions of an Event A3 handover. A
Proximal Policy Optimization (PPO) agent is then inserted at completed A3
decision opportunities to decide whether the candidate handover should be
accepted or rejected. All reported policies share the same traces, radio
processing, handover delays, radio-link-failure logic, and evaluation code.

\subsection{Project Objective and Scope}

The objective is to learn a handover policy that improves connection
reliability and reduces unnecessary mobility events without causing a large
loss in delivered rate. In particular, the project considers four outcomes:
estimated data rate, radio link failure (RLF), handover failure (HOF), and
ping-pong (PP) handover. These objectives are coupled. Triggering every valid
A3 event generally follows the strongest filtered-RSRP candidate more closely,
but it can cause unnecessary or unsuccessful handovers. Rejecting too many A3
events avoids handover interruption but can leave the UAV connected to a
deteriorating serving link. The PPO policy is therefore trained to learn an
intermediate strategy from the current radio state, timer state, and recent
handover history.

The implementation is intentionally trace-driven. It does not reproduce a
complete physical layer, packet scheduler, or RRC protocol stack. Instead, it
provides an auditable and reproducible environment for comparing decision
policies under identical radio measurements. Consequently, the reference
implementation is described throughout this report as a
\emph{3GPP-inspired A3-based simulator}, rather than as a fully conformant 5G
NR implementation.

\subsection{Trace Dataset and Preprocessing}
\label{subsec:trace-data}

\subsubsection{Dataset composition}

The current experiment uses the normalized UAV dataset stored in
\texttt{data/uav\_processed}. It contains 120 paired trajectories, where each
pair consists of one RSRP file and one SINR file. The dataset is balanced over
four UAV speeds: 30, 50, 70, and 90~km/h, with 30 trajectories at each speed.
Each trajectory contains 6000 samples collected at a fixed interval of 10~ms,
corresponding to 60~s of simulated motion. Every sample contains measurements
from 21 base stations.

The raw files in \texttt{data/UAV\_sim} are converted without changing their
measurement values. The conversion process pairs RSRP and SINR files, verifies
that their sample indices and base-station columns match, and writes the common
schema
\texttt{sample,time\_ms,bs\_0,...,bs\_20}. The loader subsequently rejects
missing pairs, non-finite values, inconsistent shapes, non-contiguous sample
indices, and sampling intervals other than 10~ms.

\subsubsection{Radio streams}

One trajectory is represented internally by three explicit streams. The first
is Layer-3-filtered RSRP, used exclusively for Event A3 detection. The second
is the native SINR trace used for spectral-efficiency and rate estimation. The
third is the SINR stream used as a radio-link-monitoring (RLM) proxy. In the
current version, the rate and RLM streams contain the same unfiltered SINR
samples, but they remain separate interfaces to allow a future RLM model to be
introduced without modifying the rate model.

Layer-3 filtering is applied independently to every base station using an
exponential recursion. The configured coefficient is $K=16$, corresponding to
a filter weight of $2^{-K/4}=1/16$. Filtering is applied only to RSRP; SINR is
not resampled or shifted. This separation ensures that A3 uses a smoothed
measurement while rate and timer events retain the native sampling instants.

\subsubsection{Available and unavailable scenario information}

The repository records the number of base stations, UAV speeds, sampling
interval, trace duration, and all RSRP/SINR measurements required by the
handover experiment. However, the normalized files do not themselves contain
the UAV altitude, base-station coordinates and heights, antenna patterns,
transmit powers, propagation-model configuration, or the route-generation
procedure. These quantities should be recovered from the original UAV
simulator configuration and reported separately before publication. The
current results are valid as trace-level comparisons, but the missing scenario
metadata limits independent reproduction of the radio traces themselves.

\subsection{Construction of the 3GPP-Inspired Handover Simulator}
\label{subsec:a3-simulator}

\subsubsection{Event A3 detection}

At each sample, the simulator identifies the non-serving base station with the
largest filtered RSRP as the candidate target. This strongest-neighbor rule is
a deterministic simulator choice; 3GPP measurement reporting can include
multiple candidate cells. The A3 entering and leaving conditions follow the
inequalities specified in TS~38.331~\cite{3gpp38331}:

\begin{align}
M_n + O_{fn}+O_{cn}-H_{\mathrm{ys}}
    &> M_s+O_{fs}+O_{cs}+O_{\mathrm{off}}, \\
M_n + O_{fn}+O_{cn}+H_{\mathrm{ys}}
    &< M_s+O_{fs}+O_{cs}+O_{\mathrm{off}},
\end{align}

where $M_n$ and $M_s$ are the filtered target- and serving-cell RSRP values.
The frequency- and cell-specific offsets are zero in the present experiment.
The hysteresis and A3 offset are both 1~dB; hence, A3 entry requires the target
RSRP to exceed the serving RSRP by more than 2~dB. State is preserved inside
the hysteresis deadband.

After A3 entry, the condition must remain valid for a time-to-trigger (TTT) of
40~ms. Leaving the event before TTT completion resets the timer. If TTT
completes, the current target is latched and a decision point is generated. The
reference Always-HO policy immediately accepts every completed event. The
action-controlled engine instead pauses protocol time until the PPO policy
chooses HANDOVER or STAY. A STAY action suppresses the same serving-target
event until A3 leaves or a different target becomes eligible, preventing the
same rejected event from being offered repeatedly.

\subsubsection{Handover preparation and execution}

An accepted handover enters a 50~ms preparation phase followed by a 40~ms
execution phase. The latched target cannot be silently changed by later
measurements. At the end of execution, the handover is considered successful
when the target-cell SINR exceeds the configured in-sync proxy. Otherwise, the
attempt is counted as an HOF and the RLF recovery process begins. These fixed
delays and the final SINR check are explicit simulation assumptions that
approximate signalling, synchronization, and target access; the simulator does
not model random access, T304, or individual RRC messages in full detail.

\subsubsection{Radio link monitoring and recovery}

RLM follows the N310/N311/T310 state progression of TS~38.331 and the relevant
measurement concepts of TS~38.133~\cite{3gpp38133}. Ten consecutive
out-of-sync indications start the 1000~ms T310 timer. Three consecutive
in-sync indications stop T310, whereas expiry of T310 declares an RLF. A
sample between the two thresholds produces neither indication. Because the
dataset does not contain hypothetical PDCCH BLER or lower-layer indication
periods, $Q_{\mathrm{out}}=-8$~dB and $Q_{\mathrm{in}}=-6$~dB are used as
explicit SINR proxies. They must not be presented as standardized 3GPP SINR
thresholds. RLF recovery is represented by a fixed 200~ms delay followed by a
suitable-cell search; complete T311 and RRC re-establishment signalling are
outside the present model.

\subsubsection{Rate and event metrics}

The trace does not provide scheduled packet throughput. The reported rate is
therefore an estimated Shannon upper bound computed from the serving-cell SINR
with a bandwidth of 10~MHz:

\begin{equation}
C[k] = B\log_2\!\left(1+10^{\mathrm{SINR}[k]/10}\right).
\end{equation}

Disconnected and interrupted intervals contribute zero delivered rate. The
relative-rate metric divides the policy's mean estimated rate by an oracle
upper bound that selects the maximum instantaneous SINR among all 21 base
stations at every sample, without handover constraints or interruption. It is
therefore a normalization reference, not a realizable handover policy.

A PP event is defined as a successful $A\rightarrow B\rightarrow A$ return
within 1~s. The PP rate is the number of PP returns divided by the number of
successful handovers. HOF rate is the number of failed executions divided by
the number of started handover executions. RLF is reported both as a count and
as events per 100~s, since an RLF can occur even when no handover is active.

\begin{table}[t]
\centering
\caption{Protocol and radio configuration used in the current experiment.}
\label{tab:protocol-configuration}
\begin{tabular}{@{}ll@{}}
\toprule
Parameter & Value \\
\midrule
Radio sampling interval & 10~ms \\
Carrier frequency & 2.1~GHz \\
Bandwidth & 10~MHz \\
Number of candidate BSs in each trace & 21 \\
Layer-3 filter coefficient $K$ & 16 \\
A3 hysteresis $H_{\mathrm{ys}}$ & 1~dB \\
A3 offset $O_{\mathrm{off}}$ & 1~dB \\
Time-to-trigger & 40~ms \\
Handover preparation time & 50~ms \\
Handover execution time & 40~ms \\
$Q_{\mathrm{out}}$ SINR proxy & $-8$~dB \\
$Q_{\mathrm{in}}$ SINR proxy & $-6$~dB \\
N310 / N311 & 10 / 3 indications \\
T310 & 1000~ms \\
Simplified RLF recovery & 200~ms \\
Ping-pong observation window & 1000~ms \\
\bottomrule
\end{tabular}
\end{table}

\subsection{Construction of the Event-Triggered PPO System}
\label{subsec:ppo-system}

\subsubsection{Decision interface and SMDP formulation}

The handover problem is formulated as a semi-Markov decision process (SMDP)
because actions occur at irregularly spaced A3 opportunities rather than at
every radio sample. A macro-step begins when A3 has remained valid for the
complete TTT and ends at the next completed A3 event or the end of the trace.
The environment advances all intermediate 10~ms protocol ticks internally.
The uncontrollable prefix before the first A3 decision is retained for
full-trace metrics but excluded from the PPO-controlled return.

The action space is binary. HANDOVER accepts the latched A3 target and starts
handover preparation, whereas STAY keeps the current serving cell and
suppresses the current event. This design does not ask PPO to discover the A3
trigger itself. PPO acts as an adaptive decision layer on top of a fixed and
auditable mobility-event generator.

\subsubsection{Causal observation design}

The observation is a normalized 19-element vector constructed only from the
current and previous samples. No future trace value is accessible to the
agent. The features are organized into four groups:

\begin{itemize}
    \item \textbf{link utility:} serving spectral efficiency, target spectral
    efficiency, and their difference;
    \item \textbf{A3 persistence and dynamics:} current, mean, and minimum A3
    margin during TTT, together with margin slope and standard deviation;
    \item \textbf{link trend and failure risk:} serving and target SINR slopes,
    serving-cell gap to $Q_{\mathrm{out}}$, target-cell gap to
    $Q_{\mathrm{in}}$, N310 progress, T310 activity, and T310 progress;
    \item \textbf{mobility history:} serving-cell residence time, time since
    the last successful handover, an indicator that the target is the previous
    serving cell, and the recent handover count.
\end{itemize}

The observation explicitly exposes both the benefit of moving to the target
and the risk of delaying the handover. Previous-cell and recent-handover
features provide direct evidence of possible ping-pong behavior, while timer
features distinguish an ordinary A3 opportunity from an urgent link-failure
situation.

\subsubsection{Reward and physical-time discounting}

The primitive reward is evaluated at every 10~ms tick. It is the normalized
delivered-rate contribution minus event penalties for a newly declared RLF
and a newly detected PP handover. The default rate scale is 100~Mbps, the RLF
penalty is 5, and the PP penalty is 1. The macro reward assigned to an action
is the sum of primitive rewards until the next decision. Handover interruption,
HOF, disconnection, and recovery are also penalized indirectly through lost
delivered rate; an HOF additionally enters RLF recovery.

Macro-steps can have different physical durations $\tau_t$. To avoid treating
a short transition and a long transition as equivalent, the discount and GAE
factors are computed as $\gamma_t=\gamma_s^{\tau_t}$ and
$\lambda_t=\lambda_s^{\tau_t}$, with per-second values
$\gamma_s=0.99$ and $\lambda_s=0.95$. PPO otherwise uses the standard clipped
probability-ratio objective of~\cite{schulman2017ppo}, a value loss, entropy
regularization, normalized advantages, gradient clipping, and optional early
stopping when the approximate KL divergence exceeds the configured target.

\subsubsection{Actor-critic architecture and training process}

The actor and critic share a multilayer perceptron with two 64-neuron hidden
layers and Tanh activations. The 19-dimensional observation is mapped to a
two-logit categorical policy head for STAY and HANDOVER and to a scalar value
head. The network contains 5635 trainable parameters. Actions are sampled from
the categorical distribution during training; deterministic evaluation uses
the action with maximum probability.

Rollouts may span multiple trajectories. Terminal masks prevent value and
advantage information from leaking between trajectory resets. Each rollout is
followed by ten epochs of minibatch optimization. The best checkpoint is
selected exclusively using mean validation return. The held-out test split is
evaluated only after checkpoint selection and never participates in model
selection.

\begin{table}[t]
\centering
\caption{PPO configuration used for the reported model.}
\label{tab:ppo-configuration}
\begin{tabular}{@{}ll@{}}
\toprule
Hyperparameter & Value \\
\midrule
Hidden layers & $64,64$ with Tanh \\
Trainable parameters & 5635 \\
Optimizer & Adam \\
Learning rate & $3\times10^{-4}$ \\
Rollout capacity & 512 macro-steps \\
Minibatch size & 64 \\
Update epochs & 10 \\
PPO clip range & 0.2 \\
Value-loss coefficient & 0.5 \\
Entropy coefficient & 0.01 \\
Maximum gradient norm & 0.5 \\
Per-second discount $\gamma_s$ & 0.99 \\
Per-second GAE factor $\lambda_s$ & 0.95 \\
Target approximate KL & 0.03 \\
Training and split seed & 17 \\
Training device & CUDA \\
\bottomrule
\end{tabular}
\end{table}

\subsection{Experimental Setup}
\label{subsec:experimental-setup}

\subsubsection{Trajectory-level data split}

Complete RSRP/SINR trajectory pairs are divided using a deterministic,
speed-stratified 70/15/15 split with seed 17. For each speed, 21 trajectories
are assigned to training, five to validation, and four to testing. The final
split therefore contains 84 training, 20 validation, and 16 test trajectories.
Splitting complete trajectories prevents samples from the same UAV route from
appearing in more than one subset. Because every test trajectory lasts 60~s,
the aggregate test duration is 960~s, with 240~s at each speed.

\subsubsection{Training execution and checkpoint selection}

The PPO run collected 100,000 SMDP macro-steps and completed 196 policy
updates. Validation was performed approximately every 5000 collected
macro-steps. Training ran on a CUDA device and required 3670.7~s, or
approximately 61.2~min, according to the recorded training log. The best
checkpoint was obtained at 45,056 macro-steps, where the validation mean
return was $-2.069\pm9.466$ over 20 trajectories. At that checkpoint, the
validation set produced an average estimated rate of 24.476~Mbps, 44 RLF
events, 73 PP events, and a handover-action fraction of 71.8\%. Later updates
did not improve validation return; therefore, the 45,056-step checkpoint was
copied to the final model and used for held-out testing.

\begin{table}[t]
\centering
\caption{Dataset split and recorded training execution.}
\label{tab:dataset-training-setup}
\begin{tabular}{@{}ll@{}}
\toprule
Item & Value \\
\midrule
Total RSRP/SINR trace pairs & 120 \\
UAV speeds & 30, 50, 70, 90~km/h \\
Trajectories per speed & 30 \\
Samples per trajectory & 6000 \\
Physical duration per trajectory & 60~s \\
Train/validation/test fraction & 70/15/15 \\
Train/validation/test size & 84/20/16 \\
Split per speed & 21/5/4 \\
Held-out test duration & 960~s \\
Target macro-steps & 100,000 \\
Completed PPO updates & 196 \\
Validation interval & approximately 5000 steps \\
Selected checkpoint & 45,056 macro-steps \\
Recorded training time & 3670.7~s (61.2~min) \\
Execution device & CUDA \\
\bottomrule
\end{tabular}
\end{table}

\subsubsection{Compared policies and fair benchmarking}

Three deterministic policies are evaluated exactly once on each test
trajectory:

\begin{itemize}
    \item \textbf{A3 Always-HO:} accepts every completed A3+TTT event. This
    policy is tick-for-tick equivalent to the audited A3 reference engine;
    \item \textbf{Always-Stay:} rejects every completed A3+TTT event. It is a
    lower-complexity control showing the consequence of suppressing all
    handovers;
    \item \textbf{PPO:} loads the validation-selected checkpoint and selects
    the maximum-probability action at every completed A3 event.
\end{itemize}

All policies use identical input traces, initial conditions, filtering,
protocol timers, interruption rules, reward parameters, and trajectory order.
The primary comparison is paired by trajectory against A3 Always-HO. Event
counts are pooled before computing rates: HOF rate is total failures divided
by total attempts, PP rate is total PP returns divided by total successful
handovers, and RLF rate is $100$ times total RLF events divided by total
simulated seconds. Mean controlled return begins at the first A3 decision,
whereas rate, RLF, PP, and HOF cover the complete trace.

\subsection{Experimental Results}
\label{subsec:experimental-results}

\subsubsection{Aggregate held-out test performance}

Table~\ref{tab:aggregate-results} reports the aggregate performance on the 16
held-out UAV trajectories. A3 Always-HO obtains the highest mean estimated rate
at 21.902~Mbps. PPO achieves 21.627~Mbps, a reduction of only 0.274~Mbps or
1.25\%. The mean relative rate decreases from 92.717\% to 91.386\%, equivalent
to 1.331 percentage points.

The small rate loss is accompanied by substantially fewer adverse mobility
events. PPO reduces RLF from 60 to 43 events (28.3\%), PP handovers from 138 to
42 (69.6\%), and HOF from 60 to 34 (43.3\%). The pooled RLF rate decreases from
6.250 to 4.479 events per 100~s. The HOF rate decreases from 12.96\% to 10.93\%,
and the PP rate decreases from 34.24\% to 15.16\%. PPO accepts 75.4\% of its
available decision actions and improves mean controlled return by 11.148,
from $-16.102$ to $-4.954$. Its connected-time fraction also increases from
96.82\% to 97.81\%.

\begin{table*}[t]
\centering
\caption{Aggregate results on the 16 held-out UAV trajectories.}
\label{tab:aggregate-results}
\resizebox{\textwidth}{!}{%
\begin{tabular}{@{}lrrrrrrrrrr@{}}
\toprule
Policy & Rate & Relative & Connected & RLF & RLF/100~s & PP & HOF & HOF rate & HO action & Return \\
 & (Mbps) & rate (\%) & time (\%) & & & & & (\%) & (\%) & \\
\midrule
A3 Always-HO & 21.902 & 92.717 & 96.821 & 60 & 6.250 & 138 & 60 & 12.959 & 100.0 & $-16.102$ \\
Always-Stay & 18.179 & 75.640 & 98.042 & 94 & 9.792 & 0 & 0 & -- & 0.0 & $-20.336$ \\
PPO & 21.627 & 91.386 & 97.808 & 43 & 4.479 & 42 & 34 & 10.932 & 75.4 & $-4.954$ \\
\bottomrule
\end{tabular}%
}
\end{table*}

Always-Stay produces no handover failure or PP event by definition because it
never starts a handover. However, it has the lowest estimated rate and the
largest number of RLF events. Its 94 RLF events are 56.7\% more than A3
Always-HO and more than twice the PPO count. This control demonstrates that the
PPO improvement cannot be reproduced by simply suppressing all A3 events.

\subsubsection{Performance by UAV speed}

Table~\ref{tab:results-by-speed} compares PPO with A3 Always-HO at each UAV
speed. The test set contains four trajectories and 240~s at every speed. PPO
maintains an estimated-rate difference below 0.6~Mbps in all four groups while
reducing RLF, PP, and HOF counts at every evaluated speed.

At 30~km/h, PPO suppresses most unnecessary returns: PP events decrease from
34 to one, while RLF decreases from seven to six. At 50~km/h, PPO halves the
number of HOFs and reduces PP events from 38 to eight. The largest RLF
improvement occurs at 70~km/h, where the count decreases from 24 to 12. At
90~km/h, the improvements are smaller: RLF decreases from 15 to 14 and HOF
from 15 to 12, although PP still decreases from 32 to 20. The 90~km/h group
therefore indicates that the learned policy remains useful at high mobility
but has less margin for suppressing failures than at the lower speeds.

\begin{table*}[t]
\centering
\caption{Held-out results by UAV speed; each speed contains four 60-s trajectories.}
\label{tab:results-by-speed}
\begin{tabular}{@{}rrlrrrrrrr@{}}
\toprule
Speed & Test time & Policy & Rate & Relative & RLF & PP & HOF & HO action & Return \\
(km/h) & (s) & & (Mbps) & rate (\%) & & & & (\%) & \\
\midrule
30 & 240 & A3 Always-HO & 23.076 & 94.777 & 7 & 34 & 7 & 100.0 & $-6.013$ \\
30 & 240 & PPO          & 22.828 & 93.208 & 6 & 1  & 4 & 70.0  &  3.338 \\
\midrule
50 & 240 & A3 Always-HO & 19.903 & 92.659 & 14 & 38 & 14 & 100.0 & $-17.062$ \\
50 & 240 & PPO          & 19.336 & 90.107 & 11 & 8  & 7  & 71.8  & $-6.152$ \\
\midrule
70 & 240 & A3 Always-HO & 21.220 & 92.504 & 24 & 34 & 24 & 100.0 & $-27.485$ \\
70 & 240 & PPO          & 21.120 & 92.081 & 12 & 13 & 11 & 73.2  & $-7.295$ \\
\midrule
90 & 240 & A3 Always-HO & 23.409 & 90.930 & 15 & 32 & 15 & 100.0 & $-13.849$ \\
90 & 240 & PPO          & 23.226 & 90.149 & 14 & 20 & 12 & 83.3  & $-9.708$ \\
\bottomrule
\end{tabular}
\end{table*}

The paired trajectory analysis supports the aggregate result. Compared with
A3 Always-HO, PPO improves controlled return on 81.25\% of test trajectories
and reduces PP events on 93.75\% of them. It produces a higher estimated rate
on only 18.75\% of trajectories, which confirms that the trained objective
primarily learns a reliability-stability trade-off rather than a rate-maximizing
policy.

\subsection{Interpretation, Limitations, and Required Follow-Up}
\label{subsec:limitations}

The present experiment provides evidence that an event-triggered PPO policy can
reject a useful subset of completed A3 events. Under the evaluated traces, the
selected policy retains approximately 98.75\% of the A3 baseline's estimated
rate while materially reducing RLF, HOF, and PP events. The improvement is
obtained without modifying A3 thresholds during an episode; adaptation occurs
through the accept-or-stay decision at each completed A3 opportunity.

The result should nevertheless be treated as a current project result rather
than a final statistical claim. Training used one random seed and one
validation-selected checkpoint, and the tables do not yet include confidence
intervals across independent training runs. The following items should be
completed before the final paper is submitted:

\begin{itemize}
    \item repeat training with multiple independent seeds and report mean,
    standard deviation, and confidence intervals;
    \item recover and document the UAV altitude, topology, antenna, transmit
    power, channel, and trajectory-generation settings used to generate
    \texttt{data/UAV\_sim};
    \item distinguish consistently between estimated Shannon rate and measured
    or scheduled throughput;
    \item retain the terms ``SINR proxy'' and ``3GPP-inspired A3 simulator''
    when describing RLM and protocol compliance;
    \item investigate sensitivity to TTT, hysteresis, A3 offset, reward weights,
    and PPO hyperparameters;
    \item add ablation studies for the timer, trend, and handover-history
    observation groups; and
    \item evaluate generalization to unseen speeds, radio scenarios, or base
    station deployments rather than only held-out trajectories from the same
    generated dataset.
\end{itemize}

In its current form, the project already provides a complete reproducible
pipeline from UAV trace conversion to A3 validation, event-triggered PPO
training, validation-only checkpoint selection, and matched held-out
benchmarking. The principal current conclusion is not that PPO universally
dominates the A3 baseline, but that it achieves a favorable and interpretable
reliability-rate trade-off in the evaluated UAV mobility scenarios.
